%------------------------------------------------------------------------------
% \PART 2
%------------------------------------------------------------------------------
\chapter[Анализ технического задания]{АНАЛИЗ ТЕХНИЧЕСКОГО ЗАДАНИЯ}

%------------------------------------------------------------------------------
% \PART 2.1 Text
%------------------------------------------------------------------------------
\section{Постановка задачи}

\hspace*{12.5 mm}В данной работе основное внимание уделяется задаче 
распознавания изображений рукописных цифр с использованием нейронных сетей, 
оптимизированных для аппаратной реализации на ПЛИС (FPGA).

Целью дипломной работы является проектирование и реализация компактной 
нейронной сети прямого распространения, способной обрабатывать изображения 
размером 28$\times$28 пикселей, представляющих рукописные цифры из 
стандартного набора данных MNIST. Основной задачей является не только 
достижение высокой точности классификации, но и обеспечение совместимости с 
аппаратной реализацией, особенно в условиях ограниченных вычислительных и 
энергетических ресурсов.

В качестве основы для нейросетевой архитектуры используется метод Learned 2D 
Separable Transform, предложенный в работе~\cite{LST2D}. Преимущество данного 
подхода заключается в использовании разделяемых по строкам и столбцам операций, 
что позволяет существенно снизить количество параметров модели по сравнению с 
традиционными полносвязными слоями. LST-архитектура даёт возможность эффективно 
использовать параллелизм, характерный для FPGA, и обеспечивает высокую скорость 
обработки без значительных потерь в точности.

Для аппаратной реализации модели выбрана микросхема XC7Z010CLG400-2, входящая в 
состав популярной отладочной платы Zybo Z7. Это решение основано на архитектуре 
Xilinx Zynq-7000, сочетающей возможности ПЛИС и встроенного процессора ARM, что 
делает его удобным для интеграции программно-аппаратных систем.

Вид сверху отладочной плата Zybo Z7 на базе Xilinx Zynq-7000
представлен на рисунке~\ref{fig:z7}~\cite{ZYBO}.

\insertfigure{z7}{pic/zybo-z7.png}{ПЛИС Zybo Z7}{9cm}

Внутренняя структура Zybo Z7 представлена на рисунке~\ref{fig:z7-in}.

\insertfigure{z7-in}{pic/zybo-in.jpg}{Структура ПЛИС Zybo Z7}{16.0cm}

На схеме продемонстрированы основные блоки SoC Zynq-7000, включая: процессорную 
систему (PS), программируемую логику (PL), AXI-интерфейсы, контроллеры памяти и 
периферийные модули. В рамках настоящего проекта будет использоваться 
взаимодействие между PL и PS через AXI-шину

Для взаимодействия с моделью и удобства проведения тестирования будет 
использоваться программная платформа PYNQ, позволяющая запускать и отлаживать 
алгоритмы с помощью Jupyter Notebook и Python API. Таким образом, будет 
реализован полный цикл: от проектирования и обучения нейросети – до её 
верификации и аппаратной имплементации.

Постановка задачи и предъявленные требования представлены в графическом 
материале ГУИР 431289.001 ПЛ.
